\chapter*{국문요약}
\addcontentsline{toc}{chapter}{국문요약}
\begin{center}\large\bfseries\thesistitlekr\end{center}

표면 마무리 공정에서 부품 전체의 균일 처리와 지정된 영역 주변의 국소 재작업은 서로 다른 요구를 가진다. 국소 폴리싱에서는 대상 영역에 맞추어 공구 이동, 접촉 힘, 작업 시간을 함께 구성해야 한다. 본 연구는 현재 장비로 미세 스크래치를 자율 판별하는 대신 사람이 처리 영역을 지정하는 설정을 사용한다. 공간 지시 이후에도 남는 동작·힘·시간 구성 문제에 대해, 규칙 기반 전략과 교시 재생에 비해 관측 조건부 학습이 어떤 수행 가치를 추가하는지 검토한다.

추적 기반 교시 장치로 공구 위치·자세, 3축 힘, RGB 영상을 시간적으로 대응하여 기록한다. 얼룩 중심 상대위치를 학습하고 얼룩 방향을 학습 기준에 정렬한다. Flow Matching 정책은 영상, 로봇 상태, 측정 힘 이력으로 위치·자세·목표 힘 시퀀스를 생성한다. 공통 실행 계층은 좌표 복원, 계획 연결, 명령 후처리를 수행하고 기존 하위 힘제어기가 접촉을 조절한다. 스킬 생성, 접촉 제어, 절차 완료 판정의 역할을 분리한다.

E1/E2는 방향 정규화 없이 평가한다. E1은 고정된 기준조건에서 규칙 기반 경로, 힘 포함 교시 재생, 힘 조건부 정책을 비교한다. 세 방법 모두 각 5회에서 접촉--부하 통과--이탈 절차를 완료했다. 교시 재생은 학습 정책보다 평균 완료시간과 시간 분산이 작았고, 학습 정책의 초과 힘 적분은 더 컸다. E2는 동작 전용 학습과 힘 관측 유무에 따른 동작--힘 학습 구성을 비교한다. 힘 관측 정책은 힘 관측을 차단한 구성보다 평균 작업자 기준 완료시간이 \BCPercent\% 짧았다. 동작 전용 구성의 한 시행은 이탈로 전환하지 못해 작업자가 중단했다. 결과는 인간 지정 국소 작업에서 학습 스킬의 실행 가능성과 구성별 시간·단계 진행·힘 사용의 차이를 보여준다. 고정 작업에서는 규칙과 교시 재생도 유효한 대안이었다. 평가는 접촉 절차 수행에 한정되며 표면 품질 향상이나 조건 변화 강건성을 입증하지 않는다.

\vspace{0.4cm}
\noindent\textbf{주요어:} 국소 로봇 폴리싱, 인간 작업 지정, 동작--힘 스킬, 모방학습, Flow Matching, 접촉 작업 평가
